\documentclass[11pt]{article}
\usepackage[margin=1in]{geometry}
\usepackage[T1]{fontenc}
\usepackage{lmodern}
\usepackage{booktabs}
\usepackage[numbers]{natbib}
\usepackage{amsmath,amssymb}
\usepackage{hyperref}
\usepackage{xcolor}
\newcommand{\pending}[1]{\textbf{\textcolor{red}{[#1]}}}

\title{Sort the Vocabulary: Prefix-Ordered Token Classes Make\\Two-Level Language-Model Heads Cheap and Faithful}
\author{Lucas Tinoco\thanks{Independent researcher. Implementation, experiments and drafting were carried out with the assistance of an AI coding agent (Claude); all results were run and checked on the author's machine. Code, extracted data and raw results will be released with this preprint.}}
\date{October 2026}

\begin{document}
\maketitle

\begin{abstract}
In small language models with large subword vocabularies, the output head is a large share of the per-token cost: 21--63\% of the weight bytes read per generated token in Gemma~3 270M, Qwen3-0.6B and Llama~3.2~1B. A two-level (class-factored) softmax cuts that cost by $\sqrt{V}$-fold, and has returned in 2025--26 in speculative-decoding drafters and drop-in head replacements; what remains open is \emph{how to form the classes}. We show that the simplest possible rule---sort the subword vocabulary by spelling and cut it into equal blocks---gives the most faithful factored head among the criteria in current use. Swapping the head of five frozen open models (Gemma~3 270M/1B, Gemma~4 E2B, Qwen3-0.6B, gpt-oss-20b; three vendors, four tokenizers), prefix-ordered classes have the lowest KL divergence to the original model in four of five (tied with semantic $k$-means on the fifth), and 27--33\% lower KL than the token ordering Google ships with its Gemma~4 drafters, at equal cost, including on Gemma~4 itself. Trained from scratch, prefix classes beat frequency, merge-order, random and semantic classes in recurrent, MLP and Transformer models, in Portuguese and English, with the gap shrinking as the vocabulary grows. A per-decision and word-start/continuation decomposition explains why: the first decision follows the generation direction, and the gain concentrates on tokens that continue a word---which also explains why alphabetical classes lost at the whole-word level (Jiang et al., 2017). On a CPU at batch size one, the factored head makes Gemma~3 270M generate 2.4$\times$ faster end to end (2.1$\times$ with AVX2 only).
\end{abstract}

\section{Introduction}

The output projection of a language model maps a hidden state $h\in\mathbb{R}^d$ to logits over a vocabulary of size $V$. Vocabularies have grown (128k in Llama~3, 152k in Qwen3, 201k in gpt-oss, 262k in Gemma) while the models people run locally have stayed small, so in small models the head is no longer a rounding error: it is 63\% of the bytes read per token in Gemma~3 270M (Table~\ref{tab:share}). At batch size one---a model on a laptop or phone, or a draft model in speculative decoding---generation is memory-bound, and those bytes are time and energy.

Class-factored softmax \citep{goodman2001classes,mikolov2011extensions} is the classical remedy: predict a class, then a token inside the class, each with a small softmax; with $C\approx\sqrt{V}$ classes of $m\approx\sqrt{V}$ tokens the head reads $(C+m)d$ instead of $Vd$ weights. It fell out of use on GPUs with large batches, where the full head is cheap, and because it costs quality. It is back: Gemma~4's drafters score 2{,}048 token clusters and expand the top 32 \citep{gemma4}; FlashHead replaces heads with balanced spherical $k$-means clusters and multi-probe retrieval \citep{flashhead}; CORAL factors a drafter head into equal groups \citep{coral}; FR-Spec and VocabTrim restrict drafter heads by frequency \citep{frspec,vocabtrim}. All of these group tokens by embedding similarity, frequency, or token id.

We ask which grouping makes the factored head most faithful, at equal cost, and find that a rule needing no data and no model---\emph{sort the vocabulary by spelling, cut into equal blocks}---wins. Our contributions:
\begin{itemize}
\item \textbf{Retrofit evidence on five frozen open models} (Section~\ref{sec:retrofit}): prefix classes give the lowest KL to the original head in four of five models (a tie in the fifth), and beat the ordering shipped with Gemma~4 by 27--33\% in KL at equal cost.
\item \textbf{From-scratch evidence} across MLP, recurrent and Transformer models, two languages and vocabularies of 8k--100k (Section~\ref{sec:scratch}), including the honest finding that the gap shrinks with vocabulary size.
\item \textbf{A mechanism} (Section~\ref{sec:why}): the class decision becomes easy because it follows the generation direction; the gain is 4--6$\times$ larger on word-continuation tokens. This reconciles our result with the opposite whole-word finding of \citet{jiang2017}.
\item \textbf{End-to-end cost} on CPU and a consumer GPU at batch one (Section~\ref{sec:cost}).
\end{itemize}

\section{Method}

\paragraph{Two-level head.} Given a partition of the vocabulary into classes $c(t)$ with members $M_c$, the head computes
\[
\log p(t\mid h) = \log \mathrm{softmax}(A h + b)_{c(t)} + \log \mathrm{softmax}\big(W_{M_{c(t)}}\, h\big)_{t},
\]
where $A\in\mathbb{R}^{C\times d}$ predicts the class and $W_{M_c}$ are the output rows of the class members. The factorization is exact: no token is ever excluded. Sampling draws a class, then a token. At inference the head reads $C\cdot d$ class weights and one block of $m\cdot d$ token weights.

\paragraph{Groupings compared.} \emph{Prefix}: sort tokens by their byte spelling, cut into equal blocks of $m$. \emph{Frequency}: sort by training frequency (the criterion of adaptive softmax and frequency-trimmed drafters). \emph{Token id / merge order}: the tokenizer's own order (CORAL-style groups). \emph{Random}. \emph{Semantic}: balanced $k$-means on the output embeddings (as in FlashHead and, by our inspection, Gemma~4). \emph{Gemma~4}: the \texttt{token\_ordering} buffer released with the Gemma~4 E4B drafter, cut into its 2{,}048 blocks of 128.

\paragraph{Retrofitting a frozen model.} We keep the model and its output rows frozen and train only $A,b$. In the \emph{self} setting the class predictor is trained to reproduce the model's own class marginals $p(c\mid h)=\sum_{t\in M_c}p_{\text{full}}(t\mid h)$ on training text (no labels, no other model), so its deviation from the full head is the pure cost of factoring. We report $\mathrm{KL}(p_{\text{full}}\,\|\,p_{\text{2-level}})$ in bits per character on held-out text; because $p(t\mid c,h)$ is exact, this equals the class-level KL. In the \emph{adapt} setting both the full head and the two-level head (class predictor and blocks) are fine-tuned on the same text with the same budget.

\section{Retrofitting frozen models}\label{sec:retrofit}

\paragraph{Setup.} Hidden states after the final norm are extracted with llama.cpp from 14 public-domain Portuguese novels (train 650--700k positions; test 714k characters, disjoint). Models: Gemma~3 270M-it and 1B-pt, Gemma~4 E2B-it (logit soft-capping 30 applied), Qwen3-0.6B, gpt-oss-20b (o200k tokenizer). Blocks of $m=\lceil\sqrt V\rceil$ (head cost cut 224--256$\times$), plus $m=128$ to match Gemma~4's ordering (120$\times$).

\begin{table}[t]
\centering
\caption{Fidelity of the factored head to the original model: $\mathrm{KL}(p_{\text{full}}\|p_{\text{2-level}})$ in bits/char on held-out text (lower is better). Prefix, semantic and the 128-block columns: mean $\pm$ s.d.\ over 3 seeds (2 for semantic on Gemma~3 270M); others: 1 seed.}
\label{tab:kl}
\small
\begin{tabular}{lccccc}
\toprule
model (vocab, classes) & prefix & semantic & frequency & random & token id \\
\midrule
Gemma 3 270M-it (262k, 512) & \textbf{0.1368}$\pm$.0001 & 0.1621$\pm$.0045 & -- & -- & -- \\
Gemma 3 1B-pt (262k, 512) & 0.0867$\pm$.0000 & \textbf{0.0854}$\pm$.0014 & 0.133 & 0.192 & 0.214 \\
Gemma 4 E2B-it (262k, 512) & \textbf{0.1838}$\pm$.0001 & 0.2738$\pm$.0024 & 0.271 & 0.329 & 0.364 \\
Qwen3-0.6B (152k, 390) & \textbf{0.0958}$\pm$.0001 & 0.1953$\pm$.0026 & 0.133 & 0.177 & 0.200 \\
gpt-oss-20b (201k, 448) & \textbf{0.0197}$\pm$.0001 & 0.0337$\pm$.0019 & 0.024 & 0.042 & 0.037 \\
\midrule
\multicolumn{6}{l}{\emph{equal cost, 2{,}048 classes of 128: prefix vs.\ the ordering shipped with Gemma 4 (3 seeds)}} \\
model & prefix & Gemma 4 & \multicolumn{3}{l}{prefix reduction} \\
Gemma 3 270M-it & \textbf{0.1158}$\pm$.0001 & 0.1593$\pm$.0003 & \multicolumn{3}{l}{27\%} \\
Gemma 3 1B-pt & \textbf{0.0724}$\pm$.0001 & 0.1086$\pm$.0001 & \multicolumn{3}{l}{33\%} \\
Gemma 4 E2B-it & \textbf{0.1656}$\pm$.0005 & 0.2414$\pm$.0005 & \multicolumn{3}{l}{31\%} \\
\bottomrule
\end{tabular}
\end{table}

\paragraph{Results.} Table~\ref{tab:kl}. Prefix is the most faithful grouping in four of five models; on Gemma~3 1B semantic $k$-means is marginally better (0.0854 vs.\ 0.0867, within one s.d.\ of its seed spread), but semantic classes are is worse than frequency on Qwen3 and Gemma~4. At equal cost prefix blocks deviate 27--33\% less than Gemma~4's own ordering, including on Gemma~4.

\paragraph{Bits per character, and a calibration caveat.} On base models the factored head costs a little: Gemma~3 1B-pt 1.561$\to$1.615 bpc (+3.5\%; 1.600 with 128-blocks), Qwen3-0.6B 1.918$\to$1.942 (+1.2\%), gpt-oss-20b 2.685$\to$2.697 (+0.5\%). On the instruction-tuned Gemmas the self-trained factored heads score \emph{below} the full head (Gemma~4 E2B-it: 2.604$\to$2.475), and the bpc ranking of groupings even inverts; the smoother class predictor partly corrects the instruction-tuned model's miscalibration on this out-of-domain text. We therefore use KL to the original model, not bpc, as the primary retrofit metric.

\paragraph{Equal-budget adaptation.} Fine-tuning both heads on the same text (Gemma~3 270M, Qwen3-0.6B): full 1.794 / 1.612 bpc; prefix 1.819 / 1.661 (+1.3\% / +3.0\%); semantic 1.842 / 1.751; frequency 1.896 / 1.724; random 1.872 / 1.723; token id 1.916 / 1.759. With 128-blocks on Gemma~3 270M: prefix +0.8\%, Gemma~4 ordering +3.0\%.

\section{Training from scratch}\label{sec:scratch}

We train small models on the same 14 novels with BPE vocabularies, the factored head with $C\approx m\approx\sqrt V$, and compare groupings at identical cost (raw runs: \texttt{results/alpha.jsonl} in the code release).

\begin{table}[t]
\centering
\caption{From scratch, test bpc (lower is better). Means over seeds where given.}
\label{tab:scratch}
\begin{tabular}{lccccc}
\toprule
model & prefix & semantic & merge order & frequency & random \\
\midrule
MLP, BPE 8k (3 seeds) & \textbf{2.208} & 2.258 & 2.240--2.248 & 2.240--2.248 & -- \\
Recurrent, BPE 8k (3 seeds) & \textbf{2.063} & 2.143 & 2.139$^1$ & 2.134$^1$ & 2.158$^2$ \\
Recurrent, English, BPE 8k & \textbf{1.762}$^1$ & -- & 1.813$^1$ & -- & -- \\
MLP, English, BPE 8k (2 seeds) & \textbf{1.889} & -- & 1.913 & 1.913 & -- \\
\bottomrule
\multicolumn{6}{l}{\footnotesize $^1$ one seed; $^2$ two seeds.}
\end{tabular}
\end{table}

\paragraph{Vocabulary size.} The advantage over frequency shrinks as the vocabulary grows: 3.5\% (8k), 2.2\% (32k), 1.1\% (100k; merge order 3.8\%, 2.3\%, 1.4\%). Larger vocabularies contain more whole-word tokens, and fewer word continuations, which is where prefix classes help (Section~\ref{sec:why}). The retrofit gaps on 152k--262k vocabularies are larger because a frozen body cannot adapt to the grouping.

\paragraph{Transformer.} A 2-layer Transformer with the prefix head reaches 2.063 bpc against 2.054 with the full 8k head (+0.45\%) at 4.2$\times$ fewer MACs per character.

\section{Why prefix classes work}\label{sec:why}

\begin{table}[t]
\centering
\caption{Decomposition of test bits/char (recurrent model, BPE 8k): class decision + within-class decision, split by tokens that start a word and tokens that continue one.}
\label{tab:split}
\begin{tabular}{lcc}
\toprule
grouping & word start (527k chars) & continuation (189k chars) \\
\midrule
prefix & 1.151 + 0.818 = \textbf{1.968} & 1.278 + 1.043 = \textbf{2.320} \\
merge order & 1.172 + 0.836 = 2.007 (+1.9\%) & 1.641 + 0.863 = 2.504 (+7.9\%) \\
semantic & 1.266 + 0.730 = 1.997 (+1.4\%) & 1.953 + 0.578 = 2.531 (+9.1\%) \\
frequency & 1.027 + 0.975 = 2.002 (+1.7\%) & 1.542 + 0.959 = 2.501 (+7.8\%) \\
second character & 1.220 + 0.773 = 1.994 (+1.3\%) & 1.468 + 0.916 = 2.384 (+2.7\%) \\
\bottomrule
\end{tabular}
\end{table}

Two hypotheses could explain the effect. \emph{Separability}---semantic classes force a hard choice between synonyms inside a class---is refuted: the within-class decision is \emph{easier} for semantic classes (0.69 vs.\ 0.88 bits/char overall). \emph{Direction} is supported: with prefix classes the first decision asks how the next piece of text begins, which the context predicts well, especially when a word is in progress. Sorting by the second character instead gives an intermediate result. Table~\ref{tab:split} localizes the gain: 4--6$\times$ larger on continuation tokens, and almost entirely in the class decision (1.28 vs.\ 1.54--1.95). With whole-word vocabularies every token starts a word, where the effect is small---consistent with \citet{jiang2017}, where alphabetical classes lost to frequency and Brown clusters.

\section{Cost}\label{sec:cost}

\begin{table}[t]
\centering
\caption{Output head as a share of the weight bytes read per generated token (Q8\_0 GGUF, tied embeddings read as the head).}
\label{tab:share}
\begin{tabular}{lccc}
\toprule
model & vocab $\times$ $d$ & head params & head share of bytes/token \\
\midrule
Gemma 3 270M & 262k $\times$ 640 & 168M of 268M & 63\% \\
Gemma 3 1B & 262k $\times$ 1152 & 302M of 1.0B & 30\% \\
Qwen3-0.6B & 152k $\times$ 1024 & 156M of 596M & 26\% \\
Llama 3.2 1B & 128k $\times$ 2048 & 263M of 1.24B & 21\% \\
SmolLM2-135M & 49k $\times$ 576 & 28M of 135M & 21\% \\
\bottomrule
\end{tabular}
\end{table}

\paragraph{GPU, head only.} On an RTX 4060 (fp16, CUDA graphs), scoring or sampling one token with the two-level head is 31--46$\times$ faster than the full head at batch one for the Llama~3.2~1B, Qwen3-0.6B and Gemma~3 270M head shapes, 20--31$\times$ at batch 8, and 2--10$\times$ at batch 512 with a naive per-row block gather.

\paragraph{CPU, end to end.} We wrote a minimal Gemma~3 runtime in Rust (Q8\_0 weights, AVX-512/AVX2, persistent thread pool) that matches llama.cpp to 0.3\% in bpc over 700 tokens, and swapped only the head. Table~\ref{tab:cpu} reports tokens/s on an idle Ryzen~7~7700 (Zen~4, 32\,MB L3, dual-channel DDR5), 512 tokens of test text scored one at a time, mean of three runs. The full head takes 51--57\% of each token's time; the prefix head (512 classes of 512) takes 1--2\%. Whole-model throughput rises 2.4$\times$ with AVX-512 and 2.1$\times$ with AVX2, at 1 or 8 threads; with 2{,}048 classes of 128 the gain is 2.4$\times$ as well. Package energy per token (RAPL, total, not idle-subtracted, because the spinning thread pool keeps the package busy) falls by the same factor (0.58 to 0.24\,J/token at 8 threads). Bits per character on these tokens: full 2.489, prefix head 2.435 (instruction-tuned model; see the calibration caveat above).

\begin{table}[t]
\centering
\caption{Gemma 3 270M (Q8\_0) end to end on CPU, batch one, tokens/s (mean of 3 runs of 512 tokens).}
\label{tab:cpu}
\begin{tabular}{lcccc}
\toprule
ISA, threads & full head & prefix head (512$\times$512) & speedup & full-head share of time \\
\midrule
AVX-512, 8 & 131.9 & 321.2 & 2.4$\times$ & 56\% \\
AVX-512, 1 & 75.5 & 183.5 & 2.4$\times$ & 57\% \\
AVX2, 8 & 119.2 & 256.0 & 2.1$\times$ & 52\% \\
AVX2, 1 & 58.2 & 120.5 & 2.1$\times$ & 52\% \\
\bottomrule
\end{tabular}
\end{table}

\section{Related work}

Class-based and hierarchical softmax \citep{goodman2001classes,morin2005hierarchical,mikolov2011extensions}; adaptive softmax \citep{grave2017adaptive}; tree-based hierarchical softmax with alphabetical, Huffman and Brown trees at the word level \citep{jiang2017}; prefix-tree lexicons for speech decoding search \citep{ney1992}. Screening and approximate heads: SVD-softmax \citep{shim2017svd}, L2S \citep{chen2019l2s}, FlashHead \citep{flashhead}. Drafter heads: Gemma~4 \citep{gemma4}, CORAL \citep{coral}, FR-Spec \citep{frspec}, VocabTrim \citep{vocabtrim}. Training-memory reduction of large heads without changing inference cost: Cut Cross-Entropy \citep{wijmans2025cut}. Semantic hierarchies for decoding: HDLM and TDLM. To our knowledge no prior work orders a subword vocabulary by spelling to form output classes.

\section{A frugal recurrent body (secondary study)}\label{sec:quartzo}

For CPU deployment we pair the head with a diagonal complex linear recurrence $h_t=(1-z_t)\,e^{i\theta}h_{t-1}+z_t c_t$ with input-only gates and a learned, data-independent rotation. We arrived at it independently, but it is the HGRU recurrence of HGRN \citep{qin2023hgrn}; our differences are a chrono-style log-uniform initialization of gate biases and stateful truncated-BPTT training. At tiny scale (2.2M parameters, 2 layers, 51 complex channels per layer) it matches a 2-layer Transformer with the same head (2.063 vs.\ 2.063 bpc) at 3.6$\times$ fewer MACs per character and an 816-byte inference state instead of a growing 128\,KB KV cache, and beats our Mamba-2-style baseline (2.078 bpc at 1.6$\times$ the MACs). Removing the rotation ($\theta=0$, so the real and imaginary halves become $2\times51$ independent real channels in the same 816 bytes) costs 0.8\%: 2.079 bpc (2 seeds: 2.0799, 2.0789) against 2.063 (3 seeds: 2.0614, 2.0665, 2.0608). At this scale the rotation helps, unlike the findings of HGRN2 and Griffin at larger scale with expanded real states. We also tried a synthetic associative-recall task (2--32 key--value pairs per line, 2M characters, 30 epochs): neither the recurrent core with or without rotation nor a same-width Transformer learned it within that budget (recurrent: 3--5\% accuracy on recalled characters, chance 3.8\%; Transformer: 5--7\%, declining with the number of pairs), so this probe is inconclusive and the known recall weakness of fixed-size states \citep{arora2023zoology,jelassi2024repeat} remains untested here.

\section{Limitations}

One text domain (19th-century public-domain literature) for all retrofit and most scratch results, plus English for replication. Retrofit models up to 21B parameters (3.6B active); scratch models are tiny. We evaluate likelihood and KL, not downstream tasks, and not speculative-decoding acceptance rates, which is how Gemma~4 uses its ordering (top-32 clusters, approximate). The factored head does not reduce the stored model size, only the bytes read per token. Gains at large batch on GPUs are modest.

\section*{Acknowledgments}
Thanks to Gabriel for many ideas along the way.

\bibliographystyle{plainnat}
\begin{thebibliography}{99}
\bibitem[Goodman(2001)]{goodman2001classes} J.~Goodman. Classes for fast maximum entropy training. ICASSP, 2001.
\bibitem[Mikolov et~al.(2011)]{mikolov2011extensions} T.~Mikolov, S.~Kombrink, L.~Burget, J.~\v{C}ernock\'y, S.~Khudanpur. Extensions of recurrent neural network language model. ICASSP, 2011.
\bibitem[Morin and Bengio(2005)]{morin2005hierarchical} F.~Morin, Y.~Bengio. Hierarchical probabilistic neural network language model. AISTATS, 2005.
\bibitem[Grave et~al.(2017)]{grave2017adaptive} E.~Grave, A.~Joulin, M.~Ciss\'e, D.~Grangier, H.~J\'egou. Efficient softmax approximation for GPUs. ICML, 2017.
\bibitem[Jiang et~al.(2017)]{jiang2017} N.~Jiang, W.~Rong, M.~Gao, Y.~Shen, Z.~Xiong. Exploration of tree-based hierarchical softmax for recurrent language models. IJCAI, 2017.
\bibitem[Ney et~al.(1992)]{ney1992} H.~Ney, R.~Haeb-Umbach, B.-H.~Tran, M.~Oerder. Improvements in beam search for 10000-word continuous speech recognition. ICASSP, 1992.
\bibitem[Shim et~al.(2017)]{shim2017svd} K.~Shim, M.~Lee, I.~Choi, Y.~Boo, W.~Sung. SVD-Softmax: fast softmax approximation on large vocabulary neural networks. NIPS, 2017.
\bibitem[Chen et~al.(2019)]{chen2019l2s} P.~H.~Chen, S.~Si, S.~Kumar, Y.~Li, C.-J.~Hsieh. Learning to screen for fast softmax inference on large vocabulary neural networks. ICLR, 2019.
\bibitem[Tranheden et~al.(2026)]{flashhead} W.~Tranheden, S.~Ahmed, D.~Dubhashi, J.~Matthiesen, H.~von Essen. FlashHead: efficient drop-in replacement for the classification head in language model inference. arXiv:2603.14591, 2026.
\bibitem[Gemma Team(2026)]{gemma4} Gemma Team. Gemma 4 technical report. arXiv:2607.02770, 2026.
\bibitem[Weng et~al.(2025)]{coral} Weng et~al. CORAL: learning consistent representations across multi-step training with lighter speculative drafter. arXiv:2502.16880, 2025.
\bibitem[Zhao et~al.(2025)]{frspec} W.~Zhao, T.~Pan, X.~Han et~al. FR-Spec: accelerating large-vocabulary language models via frequency-ranked speculative sampling. ACL, 2025.
\bibitem[VocabTrim(2025)]{vocabtrim} VocabTrim: vocabulary pruning for efficient speculative decoding in LLMs. arXiv:2506.22694, 2025.
\bibitem[Wijmans et~al.(2025)]{wijmans2025cut} E.~Wijmans, B.~Huval, A.~Hertzberg, V.~Koltun, P.~Kr\"ahenb\"uhl. Cut your losses in large-vocabulary language models. ICLR, 2025.
\bibitem[Arora et~al.(2023)]{arora2023zoology} S.~Arora et~al. Zoology: measuring and improving recall in efficient language models. arXiv:2312.04927, 2023.
\bibitem[Jelassi et~al.(2024)]{jelassi2024repeat} S.~Jelassi, D.~Brandfonbrener, S.~M.~Kakade, E.~Malach. Repeat after me: Transformers are better than state space models at copying. arXiv:2402.01032, 2024.
\bibitem[Qin et~al.(2023)]{qin2023hgrn} Z.~Qin, S.~Yang, Y.~Zhong. Hierarchically gated recurrent neural network for sequence modeling. NeurIPS, 2023.
\end{thebibliography}

\end{document}
